% Introduction draft — generated 2026-05-15, numbers from registry

\section{Introduction}
\label{sec:introduction}

Test-time compute scaling has emerged as a powerful paradigm for improving LLM-based code repair: rather than relying on a single generation, systems produce many candidate patches and select the best one through verification~\citep{snell2024scaling}. The bottleneck in this pipeline is the verifier. Full execution---running each candidate against the test suite in a sandboxed environment---is reliable but expensive, costing \$1--3 and several minutes per candidate on repository-level benchmarks~\citep{jimenez2024swebench}. Execution-free verification eliminates this cost, but existing approaches face a fundamental tradeoff between accuracy and diagnostic granularity.

On one end, trajectory-level reward models such as SWE-RM~\citep{shum2025swerm} score candidate patches as a whole, achieving strong binary accept/reject accuracy but discarding information about \emph{which} tests a patch fails. This coarse signal limits downstream applications: when a candidate is rejected, the repair system has no guidance for targeted re-repair. On the other end, Code World Models (CWM)~\citep{meta2025cwm,maimon2026selfexecution} predict execution outputs per test, preserving diagnostic granularity---but formulate the task as open-ended generation, requiring the model to reproduce exact output strings token by token.

This generation formulation creates an inherent fragility. When applied to repository-level code with long execution traces, complex return types, and verbose output formatting, CWM predictions encounter two systematic failure modes: \emph{token-budget exhaustion}, where the expected output exceeds the generation budget and exact-match necessarily fails, and \emph{string-value prediction errors}, where a single character mismatch in a floating-point number, nested data structure, or formatted string invalidates the prediction regardless of semantic correctness. On repository-level benchmarks with long execution traces, the majority of CWM predictions become unparseable, rendering the generation approach effectively unusable for real-world verification.

We argue that these failure modes are not inherent to execution simulation---they are artifacts of the generation formulation. Our key insight is that the verification task requires only a binary pass/fail verdict per test, not a reconstruction of the execution output. We reformulate execution simulation as \emph{per-test binary classification} (CLS): given a code context and a test case, the model predicts a single token---\textsc{Pass} or \textsc{Fail}---sidestepping token-budget exhaustion and string-value prediction entirely (Figure~\ref{fig:reformulation}).

This reformulation advantage is \emph{complexity-dependent}. On function-level benchmarks where execution outputs are short and simple, CWM failure modes rarely trigger, and the two formulations achieve near-equivalent accuracy: CLS attains 88.43$\pm$0.78\% and CWM 87.19$\pm$0.37\% across five random seeds (gap: 1.24 percentage points). However, as code complexity increases, the advantage amplifies along two dimensions. First, scaling model size from 4B to 8B widens the gap: CLS reaches 89.39\% versus CWM at 86.96\% (gap: 2.43pp, three seeds, $p = 0.0066$), suggesting that additional capacity benefits classification but yields diminishing returns for generation. Second, at repository-level complexity on SWE-bench Verified, the gap grows dramatically: CLS achieves 83.63\% accuracy versus CWM at 58.04\%, a 25.59 percentage-point advantage at matched 4B model scale and seed. This gap is driven by the generation failure modes identified above: in our controlled 4B experiments (\S\ref{sec:method}), 53.7--67.2\% of CWM predictions on SWE-bench Verified are unparseable due to token-budget exhaustion or string-value errors.

Crucially, this per-test formulation does not sacrifice patch-level accuracy for granularity. A TOST equivalence test confirms that per-test CLS, aggregated to patch-level verdicts, achieves statistically equivalent accuracy to a trajectory-level classifier ($p = 0.007$ at a $\pm$2pp margin), while additionally providing per-test diagnostic signals that trajectory-level verifiers discard. These diagnostics are most valuable when they are most needed: the per-test advantage concentrates in instances with multiple fail-to-pass tests ($T \geq 2$), which constitute 31\% of SWE-bench Verified instances but account for 55\% of the accuracy advantage. Furthermore, CLS and CWM exhibit complementary error patterns---CLS errors skew toward missed bugs while CWM errors skew toward false alarms---enabling an oracle ensemble accuracy of 90.63\%, pointing toward formulation diversity as a verification strategy.

We also establish that zero-shot LLM-as-judge baselines substantially underperform fine-tuned models (8B zero-shot: 64.40\%; 32B NF4 zero-shot: 71.69\%), confirming that task-specific fine-tuning remains essential for execution simulation regardless of formulation.

Our contributions are threefold:
\begin{enumerate}
    \item We identify a \textbf{complexity-dependent reformulation advantage}: classification and generation achieve equivalent accuracy on simple function-level code, but classification scales better with model size and to repository-level complexity where generation failure modes dominate (25.59pp gap at 4B scale on SWE-bench Verified). Furthermore, zero-shot LLM-as-judge baselines substantially underperform fine-tuned models (32B NF4: 71.69\% vs.\ fine-tuned 4B: 88.43\%), confirming that task-specific fine-tuning remains essential.
    \item We demonstrate that per-test classification provides \textbf{diagnostic granularity at equivalent accuracy}: per-test CLS matches trajectory-level accuracy (TOST $p = 0.007$) while preserving actionable per-test failure information.
    \item We characterize the \textbf{accuracy-utility frontier}, showing that per-test advantage concentrates in multi-fail-to-pass instances and that CLS/CWM error complementarity (oracle ensemble: 90.63\%) motivates formulation diversity for verification pipelines.
\end{enumerate}
